\begin{table}[htbp]
\centering
\caption{Border premium estimated separately by state}
\label{tab:state_specific}
\footnotesize
\begin{tabular}{lccc}
\toprule
Outcome & New Mexico & Colorado & Montana \\
\midrule
Revenue per capita & $+0.660^{***}$ & $-0.036$ & $+0.420^{*}$ \\
 & (0.218) & (0.280) & (0.253) \\
Dispensary density & $+0.595^{**}$ & $-0.337$ & $+0.398^{***}$ \\
 & (0.282) & (0.400) & (0.116) \\
Revenue per store & $+0.066$ & $+0.530^{**}$ & $-0.273$ \\
 & (0.206) & (0.239) & (0.174) \\
\midrule
Observations & 360 & 1449 & 422 \\
Border counties & 10 of 32 & 12 of 42 & 15 of 36 \\
\bottomrule
\end{tabular}
\begin{minipage}{0.98\textwidth}
\vspace{0.3em}
\footnotesize
\textit{Notes:} $\log(y_{ct}) = \alpha + \beta_s\,\text{Border}_c + \delta_t + \varepsilon_{ct}$, estimated
separately within each state (equivalently, the state-specific slice of a specification with state $\times$
quarter fixed effects). Border$_c$ = 1 if county $c$'s centroid lies within 75 miles of the nonlegal-state
line, with $\geq$1,000 nonlegal-state residents within that reach. Sample truncated to 2025Q2 or earlier.
County-clustered standard errors. Colorado's border premium is not statistically distinguishable from zero on
revenue per capita or dispensary density; Montana's is closer in magnitude to New Mexico's. This heterogeneity
is discussed directly in the text rather than only in the pooled estimate. $^{*}p<0.10$, $^{**}p<0.05$,
$^{***}p<0.01$.
\end{minipage}
\end{table}
